\section{Теоретическая часть}

\subsection{Filter-wise normalization}
Для свёрточных весов направление возмущения нормируется по каждому выходному фильтру. Это уменьшает влияние масштабной неоднозначности параметров и делает визуализацию surface более интерпретируемой.

\subsection{Квадратичные аппроксимации}
\paragraph{Полная квадратика ($q_1$):}
\[
q_1(\alpha,\beta)=c+u\alpha+v\beta+\frac{1}{2}(a\alpha^2+2b\alpha\beta+d\beta^2).
\]
Коэффициенты оцениваются методом наименьших квадратов; матрица второй производной проектируется на конус положительно определённых.

\paragraph{Диагональная квадратика ($q_2$):}
\[
q_2(\alpha,\beta)=c+u\alpha+v\beta+\frac{1}{2}(a\alpha^2+d\beta^2),\quad b=0.
\]

\paragraph{Hessian FD-квадратика ($q_3$):}
\[
q_3(\alpha,\beta)=L_0+g_1\alpha+g_2\beta+\frac{1}{2}(H_{11}\alpha^2+2H_{12}\alpha\beta+H_{22}\beta^2),
\]
где $g_1,g_2$ и $H_{ij}$ оцениваются по градиенту и конечным разностям.

\subsection{Ключевая гипотеза}
При малом радиусе $r$ поверхность локально близка к квадратичной, поэтому $q_1$ и $q_2$ должны давать меньшую ошибку, чем $q_3$.
